\begin{proofbox}
	\textbf{\textcolor{CProof}{思路提示}}

	要说明：若$f(x)$是奇函数或偶函数，则其定义域必然关于原点对称；且在定义域对称的前提下，可用$f(-x)$与$f(x)$的关系判定。

	\medskip
	\textbf{\textcolor{CProof}{正式推导}}
	\begin{description}[style=nextline, leftmargin=2.8em, labelsep=0.5em, itemsep=0.55em]
		\item[\textbf{步骤一（对称必要性）：}] 设$f(x)$是奇函数或偶函数。根据定义，对任意$x\in D$，必有$-x\in D$且满足相应的关系式。因此定义域$D$关于原点对称。
			\par\textbf{理由：} 奇偶函数的定义本身包含定义域对称的条件。

		\item[\textbf{步骤二（判定关系）：}] 在定义域对称的前提下，若$f(-x)=f(x)$恒成立，则$f(x)$为偶函数；若$f(-x)=-f(x)$恒成立，则为奇函数。
			\par\textbf{理由：} 这是奇偶函数定义的直接应用。

		\item[\textbf{步骤三（等价变形）：}] 将判定式移项得到另一种常用形式：偶函数：$f(x)-f(-x)=0$，奇函数：$f(x)+f(-x)=0$。这些形式在含参讨论时更为方便。
			\par\textbf{理由：} 等式变形的基本代数操作。
	\end{description}

	\medskip
	\textbf{\textcolor{CProof}{结论回扣}}

	\textbf{函数具有奇偶性的必要条件是定义域关于原点对称；在对称前提下，通过检验$f(-x)$与$f(x)$是否相等或相反，即可判定是偶函数还是奇函数。}
\end{proofbox}
